\chapter*{Preface}
\addcontentsline{toc}{chapter}{Preface}
This book tells a recent mathematical story as it actually unfolded: a sequence of conjectures, proofs, computer
verifications and, at times, refutations. It is written for undergraduate and graduate students in mathematics,
physics and engineering. The starting point is a theorem of openai/math, whose proof, produced by an AI and not peer
reviewed, is checked by the Lean~4 kernel: the triangular lattice minimizes the energy of every point configuration in
the plane, for a wide class of repulsive interactions. The end point is a series of open questions in quantum physics
and population genetics.

Every result is presented three times.
\begin{itemize}
\item First as a \emph{human proof}, step by step, at the level of an undergraduate course in analysis, linear algebra
and probability.
\item Then as a \emph{Lean~4} statement checked by the kernel.
\item Finally as a numerical computation with \emph{rusty-SUNDIALS} (CVODE solvers in Rust) and \emph{LeanFlow}
(data certificates checked by Lean), illustrated on public datasets.
\end{itemize}

A single rule guides the reading: every claim carries a \emph{confidence tier}, following the Elenchus method.
\begin{itemize}
\item \tierA{}: checked by the Lean kernel. This requires three checks: the axiom list is contained in
\lean{propext}, \lean{Classical.choice} and \lean{Quot.sound}; the kernel replays the proof (\lean{leanchecker}); and
a search for constructions that bypass the kernel finds nothing. The axiom list alone can be fooled; Chapter~1
explains how.
\item \tierAs{}: a variant of \tierA{} when the kernel replay has not (yet) been run for the module concerned; the
axiom list, the planted witnesses and the independent recompilation are still required.
\item \tierL{}: a published result, read and cited precisely.
\item \tierC{}: a human argument, not machine-checked.
\item \tierX{}: a floating-point computation. It guides research but proves nothing.
\end{itemize}
Read the tier before the sentence.

\paragraph{Honesty.}
\begin{itemize}
\item The results in this book were produced by AI agents (Claude) under the author's direction.
\item Each paper was audited by a separate AI agent; no human mathematician has reviewed them yet.
\item Several results depend on theorems of the openai/math release, which were also produced by an AI and have not
yet been peer reviewed.
\item Several chapters report \emph{negative results}: a refuted conjecture, a method that does not extend, a
transition that does not take place. They are as instructive as the others.
\item No conjecture is declared proved beyond its tier. In particular, the Sandier--Serfaty conjecture is \emph{not}
proved here: Chapter~2 explains exactly what was checked and what remains to be done.
\end{itemize}

\paragraph{What is new in this edition.} This is the English university edition of the French textbook
\emph{De la conjecture au noyau Lean}. The proofs, numbers, tiers, DOIs and caveats are identical. For undergraduate
and graduate readers, short \emph{University remarks}, each marked \tierL{}, state the general textbook form of a tool
used in the text (dominated convergence on a measure space, Bernstein's and Bochner's theorems, the Poisson summation
formula in $\R^d$, the theory of absorbing Markov chains) and give a reference; they add no new claim. Each of Chapters~1 to~5
ends with two additional harder exercises, titled ``(*) University exercise''; they are distinct from the
difficulty stars that some chapters already attach to their exercises. The figures are those of the French edition,
reused unchanged, so the labels inside some figures are in French; every caption is in English.

\paragraph{Outline.}
\begin{itemize}
\item Chapter~1 lays the foundations: lattices, completely monotone functions and formal verification.
\item Chapter~2 follows the thread of the formalizations, from universal optimality to the potentials of physics and
to the Sandier--Serfaty conjecture.
\item Chapter~3 presents a quantum surprise: for some interactions, the triangular lattice is no longer optimal once
the zero-point energy is taken into account.
\item Chapter~4 checks classical theorems of population genetics in Lean, up to the isothermal theorem of
evolutionary graph theory.
\item Chapter~5 gathers five use cases with data, code and certificates.
\end{itemize}
Each chapter ends with exercises.

\paragraph{Further reading.} Every cited paper has a Zenodo DOI. The code is public: AutoevolveAI,
rusty-SUNDIALS, LeanFlow, SocrateAI-Scientific-Genetic and SocrateAI-Scientific-Agora-LeanMaster.
